% BÀI TOÁN HÌNH HỌC - BAI_HINH_35 (TRÍCH ĐỀ 35)
\begin{ex}
% Câu 13
	Cho đường tròn $(O;R)$, đường kính $AB$. Gọi $D$ là trung điểm của đoạn $OB$. Từ $D$ kẻ dây $PQ$ của đường tròn $(O)$ vuông góc với đường kính $AB$. Lấy $M$ là một điểm bất kỳ trên cung nhỏ $AP$, dây $MB$ cắt dây $PQ$ tại $I$.
	\begin{enumerate}
		\item Chứng minh 4 điểm $D, I, M, A$ cùng thuộc một đường tròn.
		\item Chứng minh $QB^2 = BI \cdot BM$.
		\item Gọi $C$ là một điểm nằm trên dây $MQ$ sao cho $MP = MC$. Xác định vị trí $M$ trên cung nhỏ $AP$ để tổng $S = MA + MP$ đạt giá trị lớn nhất.
	\end{enumerate}
\end{ex}
